\section{Counting exact field roots}\label{prof:roots}

We express expected counts of field roots as integrals. Conditioning at
a root leaves a Gaussian matrix on its orthogonal complement, while the
scalar part of the integral determines the root's empirical moments.
An excluded root count will later bound the probability of an excluded
field. All integrated estimates below use the planted law, with the
original spectral restriction imposed when indicated.

The organization by a Gaussian root density and a conditional determinant
follows the Kac--Rice analysis of TAP equations in
\citet[Sections~3--4]{FanMeiMontanari2021} and its local-convexity refinement
in \citet[Lemma~4.8 and Appendix~A.5]{CelentanoFanMei2023}.
The critical shrinking scale, polynomial determinant bound, and
$nq^3$ exclusion rates are proved here.

At observation precision $s=r^2$, the planted experiment may be
gauged so that its spin is $\mathbf1$. If $S_X=\operatorname{diag}X$,
this gauge sends $J$ to $S_XJS_X$ and sends $h,y,m$ to
$S_Xh,S_Xy,S_Xm$. In the gauged coordinates
$J=W+\mathbf1\mathbf1^{\mathsf T}/n$ and
$h_s=s\mathbf1+\sqrt s\,g$, with $W$ a GOE matrix and $g$ an
independent standard Gaussian. We study the roots of
\[
 \Phi_J(y):=y+(1-q)m-Jm=h_s,\qquad
 m=\tanh y,\qquad q=\langle m^2\rangle.
\]
Write
\[
 b=1-q,\quad v=1-m^2,\quad V=\operatorname{diag}v,\quad
 M=\langle m\rangle,
\]
and, for the scalar Gaussian regression, put
\[
 \nu=q+s,\quad d=M+s,\quad a=\langle(y-d\mathbf1)m\rangle,\quad
 C_a=a/\nu,\quad j_0=q/\nu,\quad
 \alpha=(C_a-b)/(1+j_0).
\]
For fixed $y$, let $\mathsf Q_y$ be the conditional law
of $J=W+\mathbf1\mathbf1^{\mathsf T}/n$ given
\[
 Wm+\sqrt s\,g=y-d\mathbf1+bm.
\]
The field represented by $(y,J)$ is $\Phi_J(y)$. All root-counting
arguments in this section use $k=nr^3\ge\log^{10}n$ and small $r$.
\begin{lemma}[The exact counting density]\label{prof:root-count}
For every nonnegative Borel test $T(y,J)$,
\[
 \E_{\rm pl}\sum_{\substack{y:\,\Phi_J(y)=h_s\\
                         \det D\Phi_J(y)\ne0}}T(y,J)
 =\int (1+j_0)^{-1/2}G(y)\,
       \E_{\mathsf Q_y}\!\left[T(y,J)\mathcal W(y,J)\right]\,dy,
\]
where
\[
 G(y)=\prod_i \varphi_{d,\nu}(y_i)\exp\{n(1+j_0)\alpha^2/2\}
\]
and
\begin{equation}\label{prof:eq8}
\mathcal W(y,J)
=e^{-nb^2/2}|\det(I+(bI-J-2mm^{\mathsf T}/n)V)|.
\end{equation}
Here $\varphi_{d,\nu}$ is the density of a real Gaussian with mean
$d$ and variance $\nu$. Matrix restrictions are included in $T$.
\end{lemma}

\begin{proof}
In the conditioning above, $y$ is fixed and $m=\tanh y$ is its
deterministic profile. The root map is proper,
since $\Phi_J(y)-y$ is bounded for fixed $J,n$. Its critical target values
have Lebesgue measure zero, and the Gaussian field has an everywhere
positive smooth density. The area formula therefore expresses the expected
root sum as the integral of the observation density at $\Phi_J(y)$ times
$|\det D\Phi_J(y)|$; exhaustion by compact balls and monotone convergence
handle arbitrary nonnegative tests.

The derivative is
\[
 D\Phi_J(y)=I+(bI-J-2mm^T/n)V.
\]
The conditioned observation has covariance $\nu I+mm^T/n$, determinant
$\nu^n(1+j_0)$, and inverse
$\nu^{-1}I-mm^T/[n\nu^2(1+j_0)]$.
After dividing its density by $\prod_i\varphi_{d,\nu}(y_i)$, the log ratio
per site, apart from the determinant prefactor, is
\[
 -bC_a-\frac{b^2j_0}{2}
   +\frac{(C_a+bj_0)^2}{2(1+j_0)}
 =\frac{(1+j_0)\alpha^2}{2}-\frac{b^2}{2}.
\]
This proves the formula, including the conditional matrix law and all
nonnegative restrictions.
\end{proof}

\begin{lemma}[A determinant square bound]\label{prof:determinant-square}
For each fixed $D$, the conditional second moment
$\E_{\mathsf Q_y}\mathcal W(y,J)^2$ is at most $n^{C_D}$ on
$\|y\|\le n^D$, uniformly over the parameters under consideration.
Outside these balls the bound has a fixed polynomial envelope in
$\|y\|$.
\end{lemma}

\begin{proof}
The identity $\det(I+AB)=\det(I+BA)$ writes the determinant in
$\mathcal W$ as the symmetric determinant
$\det(I+\sqrt V(bI-J-2mm^{\mathsf T}/n)\sqrt V)$.
First compare the conditional matrix with an unconditioned GOE matrix
$W_0$. For $q>0$, put $u=m/\sqrt{nq}$ and $E=u^\perp$. The conditional
law leaves $P_EWP_E$ as a fresh GOE block. Its remaining row and scalar
are Gaussian; with $e=P_E\mathbf1/\sqrt n$ they have the representation
\[
 \begin{split}
 P_EJu={}&\frac{M}{\sqrt q}e+
       \frac{\sqrt q}{\sqrt n\,\nu}P_E(y-d\mathbf1)
       +\sqrt{\frac{s}{n\nu}}\,g_E,\\
 u^{\mathsf T}Ju={}&\frac{M^2}{q}
       +\frac{2(a+bq)}{2q+s}
       +\sqrt{\frac{2s}{n(2q+s)}}\,\xi,
 \end{split}
\]
where $g_E,\xi$ are independent standard Gaussians, also independent
of the $E$ block. Share that block with $W_0$. The row change, the
plant $ee^{\mathsf T}$, and the TAP spike are a perturbation of fixed
rank. Its norm is polynomial in $n,\|y\|$ and the fresh Gaussian
norms: $|M|/\sqrt q\le1$, $a^2\le q\langle(y-d)^2\rangle$, and
$\nu,2q+s\ge s$ with $s^{-1}$ polynomial on the present range.
When $q=0$, $m=0$ and the conditioning is independent of $W$, so the
same conclusion is immediate.

Here is a useful bound for this reduction. For symmetric $H$ and
rank-$j$ symmetric $N$,
\[
 |\det(H+N)|\le(1+\|N\|/\varepsilon)^j
                    |\det(H+\mathrm i\varepsilon I)|.
\]
It follows from singular-value interlacing after factoring
$H+\mathrm i\varepsilon I$: the factor
$H(H+\mathrm i\varepsilon I)^{-1}$ is a contraction and the remaining
change has rank $j$. Take $\varepsilon=n^{-5}$. Markov's polynomial
derivative inequality on $[-1/n,1/n]$, followed by integration on a
small interval around a maximizer, gives for every real polynomial
$P$ of degree at most $n$
\[
 |P(\mathrm i\varepsilon)|^2
 \le n^C\int_{-1/n}^{1/n}|P(z)|^2\,dz.
\]
Indeed $\|P'\|_\infty\le Cn^3\|P\|_\infty$ on this interval; iterating
this bound controls the imaginary displacement, and an interval of
length $cn^{-3}$ controls the maximum by the integral.
Truncate enormously large fresh entries at negligible cost by crude determinant bounds.

It remains to bound the unconditioned square. For real $z$, put
$x_i=1+z+bv_i$. The exact identity is
\[
 \E_{W_0}\det\!\left((1+z)I+bV-\sqrt V W_0\sqrt V\right)^2
 =\mathbb E_H\prod_i
 \left[(x_i+\mathrm i v_i H_1)(x_i+\mathrm i v_i H_2)
       +v_i^2(|h_1|^2+|h_2|^2)\right],
\]
where $H_1,H_2$ are independent real normals of variance $1/n$ and
$h_1,h_2$ are independent circular normals with
$\mathbb E|h_j|^2=1/n$. For completeness, represent each determinant
as the top coefficient of the exponential in one set of exterior
variables $\bar\psi,\psi$. Averaging $W_0$ adds
$-[X_{11}^2+X_{22}^2+2X_{12}X_{21}+2PQ]/(2n)$ to the exponent, where
$X_{ab}=\sum_i v_i\bar\psi_i^a\psi_i^b$,
$P=\sum_i v_i\bar\psi_i^1\bar\psi_i^2$, and
$Q=\sum_i v_i\psi_i^1\psi_i^2$. Decoupling these four quadratic terms
with the displayed Gaussian variables and extracting the site
coefficients gives the identity.

Move each $H_j$ contour to $H_j+\mathrm i b$. The factors become
$1+z+\mathrm i v_iH_j$, while the two Gaussian densities cost
$e^{nb^2}$ in modulus. This exactly cancels the factor $e^{-nb^2}$
in $\mathcal W^2$. The contour move is valid because the integrand is
an entire polynomial times a Gaussian. If $e_1,e_2$ are the eigenvalues
of the two-by-two Hermitian matrix with diagonals $H_1,H_2$ and
squared off-diagonal modulus $|h_1|^2+|h_2|^2$, the remaining modulus
is at most $\prod_{j=1}^2[(1+z)^2+e_j^2]^{n/2}$, since $0\le v_i\le1$.

More explicitly, write
$R_H^2=H_1^2+H_2^2+2(|h_1|^2+|h_2|^2)=e_1^2+e_2^2$.
The Gaussian exponent is $-nR_H^2/2$. For $R_H\le1$, the inequality
$\log(1+x)\le x-cx^2$ on a fixed bounded interval bounds the remaining
integrand by $C\exp(-cnR_H^4)$; for $R_H>1$, split off a fixed compact
annulus and then use $\log(1+x)\le x-cx$ for $x$ sufficiently large,
obtaining an integrable $Ce^{-cn\min(R_H^4,R_H^2)}$ bound.
There are six real Gaussian variables, so their normalized integral is
at most $Cn^3\int_{\mathbb R^6}e^{-cn\min(|u|^4,|u|^2)}du
\le Cn^{3/2}$. The bounded-rank and real-shift reductions multiply this
by another fixed polynomial. This proves the second-moment bound.
\end{proof}

\begin{proposition}[Scalar root diagnostics]\label{prof:scalar-diagnostics}
Fix an exclusion budget $A>0$ and an accuracy for the little-oh terms.
After choosing a sufficiently small upper bound on $r$ and then a
sufficiently large $n$, the following statements hold for
$k=nr^3\ge\log^{10}n$.

The total integral of $G(y)\,dy$ is at most polynomial in $n$. Outside
mass $n^Ce^{-Ak}$ under this scalar integral,
\[
 q\asymp_A r,\qquad |M|=O_A(q),\qquad |\alpha|=O_A(q).
\]
There are absolute $L,c>0$ such that the part with
$q\ge q_{\min}\ge Lr$ has mass at most
$n^Ce^{-cnq_{\min}^3}$. For a sufficiently small fixed exclusion
budget, the retained parameters instead satisfy
$q,M=(1\pm\epsilon)r$ and $|\alpha|\le\epsilon q$, for any requested
fixed $\epsilon>0$.

Under either budget the retained empirical moments satisfy
\begin{equation}\label{prof:eq9}
 \langle m^4\rangle=3q^2-8q^3+o(q^3),\qquad
 \langle m^6\rangle=15q^3+o(q^3),\qquad
 \langle m^3\rangle=3qM+O_A(q^3).
\end{equation}
For $l=\langle ym\rangle/q$,
$|\langle y-lm\rangle|\le C_Aq^3$, with an arbitrarily small leading
constant under a sufficiently small budget. For any fixed sufficiently
small $\eta>0$, one may also retain
\[
 \langle y^2\mathbf1_{\{|y|>q^{1/2-\eta}\}}\rangle\le C_Aq^4,
 \qquad
 \frac1n\|P_{\operatorname{span}(m,\mathbf1)^\perp}y\|^2\asymp q^3.
\]
Finally, fix $\tau>0$ and $D<\infty$. There are constants $c,C>0$
such that, for every $(t_1,t_3)$ with
$\tau\le\|(t_1,t_3)\|\le n^D$ and every $\theta\in\mathbb R$, at
least $cn$ indices satisfy
\[
 |m_i|\le C\sqrt q,\qquad |y_i-lm_i|\le Cq^{3/2},\qquad
 \operatorname{dist}\left(
 t_1\frac{m_i}{\sqrt q}+t_3\frac{y_i-lm_i}{q^{3/2}}+\theta,
 \pi\mathbb Z\right)\ge c.
\]

The determinant-weighted exclusions follow from
Lemma~\ref{prof:determinant-square} after increasing the preliminary
budget. Through Lemma~\ref{prof:root-count}, they bound the expected
number of excluded regular roots. After removing the null set of
critical target fields, the probability that any root is excluded is
at most that expected count. This is the simultaneous-root meaning of
the diagnostics.
\end{proposition}

\begin{proof}
\par\smallskip\noindent\textit{Scalar comparison and exclusion of remote parameters.}\enspace
First restrict to a fixed polynomial ball in $y$. The bound
$a^2\le q\langle(y-d)^2\rangle$ and $j_0/(1+j_0)\le1/2$ give a
Gaussian tail envelope outside it. On the ball, the logarithmic
likelihood and its parameter derivatives are polynomially bounded,
since $\nu\ge s$ and $s^{-1}$ is polynomial. Divide $q,M,a$ into a
polynomial number of bins with sufficiently fine inverse-polynomial
mesh. Replacing their empirical values by a bin representative changes
the log likelihood by $O(1)$. In the transport and variance comparisons
through the retained-range bound below, the symbols
$q,M,a,b,\nu,d,C_a,j_0,\alpha$ denote these frozen representatives;
$y$ and $m=\tanh y$ remain the variables being integrated. In the
variance comparison, $j_0^{-1}$ is used only where $C_a>1$; Cauchy--Schwarz
and the polynomial ball then also bound $q^{-1}$ polynomially.

Suppose $\alpha<1$ stays a fixed distance from one. The map
$z'_i=y_i-\alpha\tanh y_i$ has Jacobian
$\prod_i(1-\alpha v_i)$; here $\alpha$ is the frozen bin value.
Direct expansion of the two Gaussian exponents gives
\[
 \frac{G(y)\,dy}{\prod_i\varphi_{d,\nu}(z'_i)\,dz'}
 =e^{O(1)-n(b\alpha+\alpha^2/2)}
       \prod_i(1-\alpha v_i)^{-1}
 \le e^{O(1)-nH(\alpha)},
\]
where
\[
 H(u)=\frac{u^2}{2}-b[-\log(1-u)-u].
\]
The inequality uses $\log(1-uv)\ge v\log(1-u)$ for $0\le v\le1$.
In particular $H(u)\ge qu^2/6$ for $u\le q$, with a cubic or stronger
penalty on the negative side when $q$ is small.

The remaining branch is $C_a>C_0:=1+j_0q$, equivalently $\alpha>q$.
Compare there with the product Gaussian on $y$ having variance
$\nu C_a^2/j_0$. Cauchy--Schwarz on $a$ gives the remaining rate
$f(C_a)/2$, where
\[
 f(C_a)=C_a^2/j_0-1-\log(C_a^2/j_0)
               -(C_a-b)^2/(1+j_0).
\]
The bounds
$f(C_0)\ge-2q^3/3$, $f'(C_0)\ge2q/(1+q)$, and $f''\ge1$
exclude $\alpha>q+C_AR_*^{3/2}$ at budget $AnR_*^3$, where
$R_*:=\max(r,q)$. The $H$ comparison excludes
$\alpha<-C_AR_*$. On the retained range
\[
 -C_AR_*\le\alpha\le q+C_AR_*^{3/2},\qquad
 H(\alpha)\ge-\delta R_*^3,
\]
where $\delta$ can be made arbitrarily small by decreasing the upper
scale.

The frozen-symbol convention for those two comparisons ends here.
We also need a fixed negative rate away from $q=0$ as $s\downarrow0$.
Fix $q_0>0$. For a probability law $P$ with finite second moment and $\E_P\tanh^2Y\ge q_0$, define
\[
 q=\E_P\tanh^2Y,\quad M=\E_P\tanh Y,\quad d=M+s,\quad \nu=q+s,
\]
\[
 a=\E_P[(Y-d)\tanh Y],\quad b=1-q,\quad j_0=q/\nu,\quad
 \alpha=(a/\nu-b)/(1+j_0).
\]
Put $\psi_s(P)=(1+j_0)\alpha^2/2$ and
\[
 \mathscr R_s(P)=\psi_s(P)
       -\operatorname{KL}\bigl(P\Vert N(d,\nu)\bigr).
\]
The endpoint scalar inequality of Lemma~\ref{prof:endpoint-scalar},
proved independently below, gives $\mathscr R_s\le0$, with equality
only for the indicated Gaussian. At $s=0$, a consistent equality law
with $q>0$ would be $N(M,q)$. Its mean equation
$M=\E_{N(M,q)}\tanh Y$ forces $M=0$, since the difference of the two
sides is strictly increasing and vanishes at zero. But then
$\E_{N(0,q)}\tanh^2Y<\E Y^2=q$, a contradiction.

For completeness, this pointwise exclusion has the uniformity needed
by the integral. Fix a bound on $\E_PY^2$. The resulting set of laws
with $q\ge q_0$ is weakly compact. Its first moments are
uniformly integrable, so the average of $y\tanh y$ and hence $\psi_s$
are continuous there; $\nu\ge q_0$ keeps the denominators separated
from zero. Relative entropy is lower semicontinuous, making
$\mathscr R_s$ upper semicontinuous. Its maximum at $s=0,q\ge q_0$
is therefore strictly negative, and remains so for sufficiently small
$s$. The second-moment bound also makes replacement of nearby Gaussian
parameters uniformly $e^{o(n)}$ in the product density. A finite cover
by neighborhoods on which one bounded continuous entropy test has
this slack bounds the scalar integral by $e^{-cn}$. This is the
compact-set integral bound in the proof of
Lemma~\ref{prof:endpoint-scalar}, with its smoothing variance set to zero
and $\nu=q+s\ge q_0$; that proof records the finite-cover argument of
\cite[Lemma~5.4]{AcceptedGap}.
The Gaussian envelope $G(y)\le e^{Cn-c\sum_i y_i^2}$ on $q\ge q_0$
removes the second-moment restriction. This proves the fixed-$q_0$
exclusion; the cubic rate for shrinking $q$ is derived next.

\par\smallskip\noindent\textit{Empirical moments at the cubic scale.}\enspace
The transport comparison handles the negative and moderate scalar
parameters, and variance inflation excludes the remaining positive
range. On the retained range we can now impose product-Gaussian
moment diagnostics and use root consistency to locate the overlap scale.

Return to the root integral. For the shifted-coordinate reference of
the bin containing $y$, write $\alpha_{\rm bin},d_{\rm bin},\nu_{\rm bin}$
for its frozen values and put
\[
 z'_{{\rm bin},i}=y_i-\alpha_{\rm bin}\tanh y_i,\qquad
 Z_i=(z'_{{\rm bin},i}-d_{\rm bin})/\sqrt{\nu_{\rm bin}}.
\]
The rate parameters in this product-reference estimate also use the
frozen bin values. For any fixed $A'>0$, exponential Markov
bounds give the following outside product probability
$e^{-A'nR_*^3}$: the second moment on $|Z_i|>R_*^{-\eta}$ is at most
$C A'R_*^3$, and clipped moments through any prescribed fixed degree
differ from their Gaussian values by $O(\sqrt{A'}R_*^{3/2})$.
Choose $\eta$ sufficiently small for that degree. These estimates use
a fixed exponential tilt for the tail square and a tilt of order
$\sqrt{A'}R_*^{3/2}$ for each clipped polynomial; their constants stay
bounded as $A'\downarrow0$.

The scalar change of measure costs at most $e^{\delta nR_*^3}$ on a
retained bin. Hence a reference exception of size $e^{-A'nR_*^3}$
contributes at most $e^{-(A'-\delta)nR_*^3}$ to the scalar integral,
up to the polynomial bin factor. For a small target budget choose
$A'$ proportional to that budget and then make $\delta<A'/2$ by
shrinking the upper scale. This preserves the small
$O(\sqrt{A'})$ moment errors needed below.

For the consistency and residual conclusions below, the symbols
$q,M,a,b,\nu,d,C_a,j_0,\alpha$ again denote the exact empirical quantities
defined at the start of the section. Expand the exact identity
$m=\tanh(z'_{\rm bin}+\alpha_{\rm bin}m)$ on the clipped region and use
the tail second moment on its complement. Replacing the bin parameters
in the resulting expressions by the empirical quantities is included
in the displayed errors, using the initially chosen sufficiently fine
mesh. The first consistency estimate is
\[
 q=\nu+d^2+O_A(R_*)(\nu+d^2).
\]
It implies $s+d^2\le C_AR_*q$. If $q\le r$, the inequality
$r^2=s\le C_Arq$ forces $q\ge c_Ar$; otherwise this is automatic.
Thus $R_*=O_A(q)$ and $|d|=O_A(q)$. On this narrower range the
retained $H$ is nonnegative. Keeping the next terms gives
\[
 0=s+d^2+2\alpha q-2q^2+o_A(q^2),\qquad
 M=d+\alpha d-qd+O_A(q^2).
\]
Apart from an $o(q^2)$ error, the last constant tends to zero with the
small budget. The $H$ and $f$ penalties then force $|\alpha|\ll q$.
Since $d=M+s$, division of the two consistency equations by $q^2$
gives
\[
 \frac{s}{q^2}=\frac dq+o_{\rm small}(1),\qquad
 \frac{s}{q^2}+\left(\frac dq\right)^2=2+o_{\rm small}(1).
\]
The first ratio is nonnegative. These two relations therefore force
$s/q^2$ and $d/q$ close to one, proving the small-budget location.
In particular, bins with $q\gg r$ have a fixed negative rate
$cnq^3$. This tail also pays a prescribed budget $Ak$ in the
large-budget calculation: choose $q_{\min}/r$ so that
$c(q_{\min}/r)^3>A$, with slack for the polynomial bins.

Finally the total scalar integral is polynomial. The exceptional bins
just considered have strict negative rates. On the retained bins the
reference product integral is one and $H\ge0$; indeed
$H(q)\ge q^3/6$ and this sign persists up to
$q+O_A(q^{3/2})$. Summing over the polynomial grid proves the claim.

\par\smallskip\noindent\textit{Fourth and sixth moments and the cubic residual.}\enspace
The scale comparison is not yet enough for the matrix estimates: they
also use the first correction to the fourth moment. We therefore keep
the cubic-order terms in the same expansion. The cubic residual will
require a separate cancellation at small exclusion budget.

The same expansion has linear coefficient $1+\alpha+O(q^2)$ and
cubic coefficient $-1/3+O(q)$. The consistency equation for $q$ gives
\[
 3q^2=3\nu^2+6\nu d^2+12\alpha\nu^2-12\nu^3+o(q^3),
\]
whereas direct expansion of the fourth moment gives
\[
 \langle m^4\rangle
 =3\nu^2+6\nu d^2+12\alpha\nu^2-20\nu^3+o(q^3).
\]
Since $\nu=q+O(q^2)$, their difference is $-8q^3+o(q^3)$.
Likewise $\langle m^6\rangle=15\nu^3+o(q^3)$ and
$\langle m^3\rangle=3\nu d+O_A(q^3)=3qM+O_A(q^3)$.
The clipped empirical errors are $o(q^3)$ at this scale even for a
fixed large budget. On the complement, boundedness of $\tanh$ and the
original-scale tail second moment $O_A(q^4)$ control the remainders.

For the residual, use the identity
\[
 \langle y-lm\rangle
 =\langle y-m\rangle-
       \frac{\langle(y-m)m\rangle}{q}M.
\]
Expand it in $z'_{\rm bin}$ through degree six, including the occurrences of
$M$ and $q$ in this expression. Replacing empirical moments by their
Gaussian values costs $O(\sqrt{A'}q^3)+o(q^3)$. The order-$q^2$
terms cancel: the cubic Gaussian moment starts with $3\nu d$ and the
fourth with $3\nu^2$. After division by $q^3$, the remaining leading
coefficient is continuous in the bounded ratios
$\nu/q,d/q,\alpha/q$, with $\nu/q$ bounded below.

At the small-budget limit these ratios approach $1,1,0$. The leading
coefficient vanishes there by an exact reflection identity. If
$Y\sim N(\nu,\nu)$, its density obeys
$\varphi(-y)=e^{-2y}\varphi(y)$, so pairing $y$ and $-y$ gives
\[
 \E F(Y)=\E[F(Y)\tanh Y]
 \quad\text{for every integrable odd }F.
\]
With $F(y)=y$ and $F(y)=\tanh y$, this gives
$\E Y=\E[Y\tanh Y]$ and $\E\tanh Y=\E\tanh^2Y$.
The integrated residual is therefore exactly zero at
$\alpha=0,d=\nu$. Continuity and the empirical error prove the stated
small leading constant.

\par\smallskip\noindent\textit{Nonlattice and spanning diagnostics.}\enspace
On bounded standardized inputs, the normalized profiles
$m/\sqrt q$ and $(y-lm)/q^{3/2}$ converge to the nonsingular pair
$z$ and $z^3/3-z$. Their Gram matrix on a fixed bounded input set is
consequently bounded above and below. Put $w=y-lm$.
Then $\langle wm\rangle=0$, $\langle w^2\rangle\asymp q^3$, and
$1-M^2/q=1-O_A(q)$ is bounded below. The normalized squared norm of
the projection of $w$ onto $\operatorname{span}(m,\mathbf1)$ is
$\langle w\rangle^2/(1-M^2/q)=O_A(q^6)$. Removing both directions
therefore leaves squared norm comparable to $nq^3$, as asserted.

For phase dispersion we need derivative control of the actual profiles,
because the frequencies may grow with $n$. In a frozen reference bin,
write $Y(z)$ for the inverse of
\[
 Y-\alpha\tanh Y=d+\sqrt\nu\,z,
 \qquad
 Y'(z)=\frac{\sqrt\nu}{1-\alpha\operatorname{sech}^2Y(z)}.
\]
The retained parameters satisfy $d=O_A(q)$, $\nu=q+O_A(q^2)$,
$\alpha=O_A(q)$, and
$l=(a+dM)/q=1+q+O_A(q^2)$. The last expansion follows by expanding
$ym=m^2+m^4/3+m^6/5+\cdots$ on the clipped region and using
\eqref{prof:eq9} and the tail-square bound. On every fixed bounded
$z$ interval, the functions
\[
 f_q(z)=\frac{\tanh Y(z)}{\sqrt q},\qquad
 g_q(z)=\frac{Y(z)-l\tanh Y(z)}{q^{3/2}}
\]
therefore satisfy
$f_q(z)=z+O_A(\sqrt q)$ and
$g_q(z)=z^3/3-z+O_A(\sqrt q)$ in $C^1$. For example,
\[
 g_q'(z)=
 \frac{(1-l\operatorname{sech}^2Y(z))\sqrt\nu}
 {q^{3/2}(1-\alpha\operatorname{sech}^2Y(z))},
\]
which converges uniformly to $z^2-1$ on that interval.

Put $R=\|(t_1,t_3)\|$ and $v=(t_1,t_3)/R$. Each polynomial
$v_1z+v_3(z^3/3-z)$, with $\|v\|=1$, has an interval of fixed
positive length on which its derivative is bounded above and away
from zero in absolute value. Compactness in $v$ permits a finite
choice of such intervals and uniform constants. The $C^1$ bounds
give the same derivative bounds for the actual function
$v_1f_q+v_3g_q$ when the upper scale is small. On its chosen interval
this function is monotone. For all sufficiently large $R$, counting
periods of $R(v_1f_q+v_3g_q)+\theta$ shows that a fixed positive
fraction of the interval stays a fixed distance from $\pi\mathbb Z$,
uniformly in $\theta$. For $\tau\le R$ in a fixed bounded range,
compactness in $R,v,\theta\bmod\pi$ gives the same conclusion with
possibly smaller constants. Gaussian density is bounded below on
the bounded intervals used. No approximation error is multiplied
by the growing frequency in this argument.

Use a smooth periodic test supported a little inside this separated
phase set and a smooth cutoff inside the bounded input interval.
For each fixed choice, product concentration gives failure $e^{-cn}$.
All parameter derivatives on the retained bins are polynomially
bounded, since $q^{-1}$ is polynomially bounded. The parameter $l$
is determined by the existing bins through $(a+dM)/q$. Refine their
mesh, and mesh frequencies and phase, at inverse-polynomial accuracy
chosen after $D$, so that every test changes by less than the fixed
separation slack even for $R\le n^D$. There are only polynomially
many tests. Their union retains failure $e^{-c'n}$, and the scalar
change of measure preserves that rate after the upper scale is
decreased. The same bounded-input concentration gives the Gram
estimates used above.
\end{proof}
